\section{Objectives}

The primary goal of this project is to evaluate the impact of architectural evolution—ranging from shallow architectures to attention mechanisms on the classification accuracy of the UCF101 action recognition benchmark.

\begin{itemize}
    \item \textbf{Quantify the Performance Floor}: Establish a baseline by training a scratch-built 4-layer CNN to measure the limits of minimal depth and lack of pre-trained spatial knowledge.
    \item \textbf{Test Spatial Transferability}: Assess the effectiveness of ImageNet-pretrained features (via ResNet18) in identifying actions that are strongly characterized by their environmental or object-based context.
    \item \textbf{Resolve Temporal Collapse}: Integrate recurrent modules (LSTM and BiLSTM) to bridge the gap between static frame-averaging and chronological sequence modeling, specifically for motion-heavy classes.
    \item \textbf{Benchmark Recurrence vs. Attention}: Implement and compare the performance of Transformer-based architectures against recurrent models to evaluate the advantages of self-attention in capturing long-range dependencies.
    \item \textbf{Class-specific Behavioral Mapping}: Generate a final performance matrix to identify which architectural paradigm is most effective for different categories of human actions (e.g., rhythmic, static, or high-speed movements).
\end{itemize}